% Standalone note. From the repository root, compile with:
% latexmk -pdf -outdir=.vscode/latex-build/notes notes/iid_cooperative_characterization.tex
\documentclass[11pt]{article}

\usepackage[T1]{fontenc}
\usepackage{lmodern}
\usepackage[margin=1in]{geometry}
\usepackage{amsmath,amssymb,amsthm,mathtools}
\usepackage{microtype}
\usepackage{enumitem}
\usepackage[round,authoryear]{natbib}
\usepackage{hyperref}

\newcommand{\X}{\mathcal X}
\newcommand{\C}{\mathcal C}
\newcommand{\V}{\mathcal V}
\newcommand{\Hclass}{\mathcal H}
\newcommand{\A}{\mathcal A}
\newcommand{\E}{\mathbb E}
\newcommand{\ind}{\mathbf 1}
\DeclareMathOperator{\Unif}{Unif}

\newtheorem{theorem}{Theorem}
\newtheorem{lemma}[theorem]{Lemma}
\newtheorem{proposition}[theorem]{Proposition}

\setlist[enumerate]{itemsep=3pt,topsep=5pt}
\setlength{\parskip}{3pt}
\allowdisplaybreaks[1]

\title{A characterization of cooperative debate with iid instances}
\author{}
\date{}

\begin{document}
\maketitle

We want to characterize when a fixed verifier class permits learning from
independent and identically distributed (iid) instances. The difficulty is
that the true label is encoded by a whole verifier table, while each round
reveals only one entry. The induced hypothesis class alone does not describe
how much these observations tell us. We must also keep track of the
constraints that the verifier class imposes between tables at different
instances.

We give a necessary-and-sufficient condition using finite restrictions of
the verifier class. On a tuple of instances, reveal one random verifier entry
at every position except a uniformly chosen held-out position. The condition
is that the worst-case error in predicting the true label at that position
tends to zero as the tuple grows. Each finite problem is a linear program.
This is an adaptation of the transductive approach to learnability; see
\citet{asilis2024transductive} for the standard supervised framework. We prove
the adaptation to verifier feedback here.

\section{The model and the learning objective}\label{sec:model}

Let $\C$ be a finite certificate set with $|\C|=k\ge1$. We consider one pair
of certificates per debate, so the query set is $\A=\C^2$ and its size is
$K=k^2$. A verifier is a function $v:\X\times\A\to\{0,1\}$. It is valid
if, for every $x$, there is a unique bit $b$ satisfying
\[
\exists p\in\C\ \forall r\in\C:\quad v(x,p,r)=b.
\]
Write $h_v(x)$ for this bit. Thus the table at $x$ has an all-$h_v(x)$ row
and no all-$(1-h_v(x))$ row.

Fix any nonempty class $\V$ of valid verifiers, and let
$\Hclass=\{h_v:v\in\V\}$ be its induced hypothesis class. We impose no
closure condition on $\V$ and do not assume that it contains all valid
verifiers for the hypotheses in $\Hclass$.

A verifier $v^*\in\V$ and a distribution $D$ on $\X$ are fixed before
learning begins. At round $t$, the cooperative learner receives
$X_t\sim D$, independently of the past. It predicts $\hat Y_t$, chooses
$A_t\in\A$, and observes only $Z_t=v^*(X_t,A_t)$. Both the prediction and
the query may depend on the history, $X_t$, and the learner's randomness.
The true label $h_{v^*}(X_t)$ is not revealed. The learner cannot make extra
queries at the current instance or at previously observed instances.
Its number of prediction mistakes is
\[
M_T=\sum_{t=1}^T\ind\{\hat Y_t\ne h_{v^*}(X_t)\}.
\]

We call $\V$ \emph{uniformly iid learnable} if there is one learner and a
bound $R(T)=o(T)$ such that $\E[M_T]\le R(T)$ for every $D$, every
$v^*\in\V$, and every $T$. The bound may depend on $\V$ and $k$, but not
on $D$ or $v^*$. This is stronger than requiring sublinear mistakes separately
for each fixed pair $(D,v^*)$, with no common rate.

We will also use passive training observations of the form
\[
(X,A,Z),\qquad X\sim D,\quad A\sim\Unif(\A),\quad Z=v^*(X,A),
\]
where $X$ and $A$ are independent. A passive learner receives $m$ iid
observations and outputs a predictor for a fresh instance. Its predictions
may be randomized, and its output need not belong to $\Hclass$. Define
$r_m(\V)$ to be the infimum, over all such learners, of their worst-case
expected prediction error, where the worst case ranges over $D$ and $v^*$.
The expectation includes training data, the fresh instance, and learner
randomness. We allow $m=0$ in this definition.

Appendix~\ref{app:passive} proves that $r_m(\V)\to0$ is equivalent to
uniform iid learnability in the original cooperative protocol. The reduction
uses that $K$ is finite and that predictions and queries can be chosen
separately. It makes no computational efficiency claim.

\section{A finite prediction problem}\label{sec:finite}

The finite problem retains every relation between verifier tables that is
visible on a given tuple. Fix $S=(x_1,\ldots,x_n)\in\X^n$, allowing
repetitions, and define
\[
W_S=\left\{
  \bigl(v(x_i,a)\bigr)_{i\in[n],\,a\in\A}:v\in\V
\right\}.
\]
The elements of $W_S$ are $n\times K$ binary arrays. The entries in row $i$
form the complete verifier table at $x_i$. Let $b_i(w)$ be the unique label
validated by that table. If two verifiers induce the same array, their
validated labels agree at every position, so $b_i(w)$ is well defined.

Consider the following game.
\begin{enumerate}
\item The learner knows $S$ and $W_S$. An adversary chooses $w\in W_S$,
before the random choices below, without revealing it.
\item Choose $I\sim\Unif([n])$ and reveal $I$.
\item For each $j\ne I$, independently choose $A_j\sim\Unif(\A)$ and
reveal $(A_j,w_{j,A_j})$.
\item The learner predicts $b_I(w)$.
\end{enumerate}
Every occurrence of an instance has its own independent query. In
particular, equal values $x_i=x_j$ at distinct indices do not share their
query randomization.

Let $\tau(S)$ be the minimax error probability in this game, allowing
randomized predictions, and set
\[
\tau_n(\V)=\sup_{S\in\X^n}\tau(S).
\]
The game is finite, so its optimum is attained and can be expressed as a
linear program. Its value depends only on $W_S$, including the labels
determined by its rows. Always $0\le\tau(S)\le1/2$, since a fair random
prediction has error $1/2$.

For measurability, assume that every verifier is measurable and, for every
$n$ and every binary array $w$, the set
$\{S\in\X^n:w\in W_S\}$ is measurable. These conditions are automatic
on a countable domain with the discrete sigma algebra. They ensure that
choosing a fixed optimal strategy for each finite pattern family gives a
measurable prediction rule. We work under these conditions throughout.

\section{The characterization}\label{sec:characterization}

The learner for the positive direction is simple to describe. Given passive
training observations and a test instance, append the test instance to the
training tuple, randomly permute all positions, and use the optimal rule for
the finite problem in Section~\ref{sec:finite}. The test position has no
revealed verifier entry. The other positions retain their observed query
and response.

\begin{theorem}\label{thm:characterization}
The verifier class $\V$ is uniformly iid learnable if and only if
$\tau_n(\V)\to0$ as $n\to\infty$.
\end{theorem}

We prove the theorem by comparing the finite problem with passive learning.
The first direction uses the learner just described.

\begin{lemma}\label{lem:batch-from-finite}
For every $n\ge1$, $r_{n-1}(\V)\le\tau_n(\V)$.
\end{lemma}
\begin{proof}
Fix $D$ and $v^*$. Take $n-1$ iid passive training observations and an
independent test instance. Append the test instance and randomly permute
all $n$ positions. Let $S$ be the resulting ordered tuple, and let $I$ be
the position occupied by the test instance.

By exchangeability, $I$ is uniform on $[n]$ and independent of $S$. The
queries at the other indices are independent and uniform in $\A$,
independently of $S$ and $I$. For each fixed $S$, the true verifier induces
an array $w\in W_S$. Conditional on $S$, the learner therefore faces
exactly the finite prediction problem, with this array as the adversary's
choice. Its conditional error is at most $\tau(S)$, and hence its
unconditional error is at most $\tau_n(\V)$.

The procedure defines a prediction rule for every test instance using the
same training data. Since the bound holds for every $D$ and $v^*$, the
claimed inequality follows.
\end{proof}

For the reverse direction, we use an iid learner on the uniform distribution
over the tuple. The finite game supplies only one query at each available
index. We can still simulate an iid sample unless we draw an index twice or
draw the held-out index.

\begin{lemma}\label{lem:finite-from-batch}
For every $n\ge1$ and $m\ge1$,
\[
\tau_n(\V)\le r_m(\V)+\frac{m(m+1)}{2n}.
\]
\end{lemma}
\begin{proof}
Fix a passive learner whose worst-case expected error after $m$ observations
is at most $r$, and fix a tuple $S=(x_1,\ldots,x_n)$. Let $D_S$ be the
distribution obtained by choosing an index uniformly from $[n]$ and
returning the instance at that index. Repeated values in $S$ contribute
their multiplicities to $D_S$.

In the finite game, draw indices $J_1,\ldots,J_m$ independently and
uniformly from $[n]$, independently of $I$ and all revealed queries. If
these indices are distinct and none equals $I$, feed their revealed
observations to the passive learner and use its prediction at $x_I$.
Otherwise, make an arbitrary prediction.

To bound the error, couple this procedure to an ideal experiment in which
the passive learner receives $m$ iid observations from $D_S$ and is tested
at an independent instance from $D_S$. Use $J_1,\ldots,J_m$ for its
training indices and $I$ for its test index. On the event that the training
indices are distinct and avoid $I$, their available queries are independent
and uniform, so they can be used as the ideal experiment's queries. On the
complementary event, generate fresh independent queries for the ideal
experiment. Conditional on every index sequence, the ideal queries have
the required law. Thus the two procedures agree whenever the simulation
succeeds.

By the union bound, its failure probability is at most
\[
\Pr\bigl(\exists\ell:J_\ell=I
  \text{ or }\exists\ell<q:J_\ell=J_q\bigr)
\le\frac mn+\frac{\binom m2}{n}
=\frac{m(m+1)}{2n}.
\]
For each $w\in W_S$, choose a verifier realizing $w$. The ideal experiment
has expected error at most $r$ for that verifier, so the finite-game rule
has error at most $r+m(m+1)/(2n)$ for every $w$. Taking the supremum over
$S$ and then the infimum over passive learners proves the claim. An optimal
passive learner need not exist for this argument.
\end{proof}

\begin{proof}[Proof of Theorem~\ref{thm:characterization}]
If $\tau_n(\V)\to0$, Lemma~\ref{lem:batch-from-finite} gives
$r_m(\V)\to0$. Conversely, if $r_m(\V)\to0$, first choose $m$ to make
$r_m(\V)$ small, then let $n$ grow in Lemma~\ref{lem:finite-from-batch}.
This gives $\tau_n(\V)\to0$. Lemma~\ref{lem:passive-equivalence} in
Appendix~\ref{app:passive} identifies $r_m(\V)\to0$ with uniform iid
learnability in the cooperative protocol.
\end{proof}

The positive direction gives an explicit online guarantee. At round $t$,
use the finite-game rule on the previous $t-1$ passive observations and
the current instance, and then query the current verifier table uniformly.
The past observations remain iid regardless of the predictions. Therefore
\[
\E[M_T]\le\sum_{t=1}^T\tau_t(\V).
\]
If $\tau_t(\V)\to0$, the right-hand side is $o(T)$.

Randomization of the prediction is not needed for the vanishing-error
criterion. Let $p(s,x)$ be a rule's probability of predicting $1$, given
training data $s$ and test instance $x$, including all its internal
randomization. Replace the rule by $g_s(x)=\ind\{p(s,x)\ge1/2\}$. For
either true bit $b$,
\[
\ind\{g_s(x)\ne b\}
\le 2\bigl(p(s,x)\ind\{b=0\}
 +(1-p(s,x))\ind\{b=1\}\bigr).
\]
Thus a deterministic output function costs at most a factor of two in the
expected error. It may be outside $\Hclass$. This rounding leaves the
uniform query choices unchanged.

\section{Computing the finite parameter}\label{sec:parameter}

The finite-game definition can be written without quantifying over a
learner. Choose a prior $\pi$ on $W_S$. For each possible observation
$o=(i,a,z)$, the index $i$ is held out, while
$a\in\A^{[n]\setminus\{i\}}$ and
$z\in\{0,1\}^{[n]\setminus\{i\}}$ list the revealed queries and
responses. For $b\in\{0,1\}$, define the prior mass consistent with this
observation and target bit by
\[
q_b^\pi(i,a,z)=
\sum_{\substack{w\in W_S:\ b_i(w)=b,\\
                   w_{j,a_j}=z_j\ \text{for all }j\ne i}}\pi(w).
\]

\begin{proposition}\label{prop:dual}
For every finite tuple $S$ of length $n$,
\begin{equation}\label{eq:dual}
\tau(S)=\max_{\pi\in\Delta(W_S)}
\frac{1}{nK^{n-1}}
\sum_{i=1}^n\ \sum_{a\in\A^{[n]\setminus\{i\}}}
\ \sum_{z\in\{0,1\}^{[n]\setminus\{i\}}}
\min\{q_0^\pi(i,a,z),q_1^\pi(i,a,z)\}.
\end{equation}
\end{proposition}
\begin{proof}
Fix a prior $\pi$. The probability of observing $(i,a,z)$ is
$(q_0^\pi(i,a,z)+q_1^\pi(i,a,z))/(nK^{n-1})$. When this probability is
positive, the two posterior target probabilities are the corresponding
$q_b^\pi(i,a,z)$ divided by their sum. The optimal prediction under this
prior chooses a bit of larger posterior probability. Its contribution to
the total error is therefore
$\min\{q_0^\pi(i,a,z),q_1^\pi(i,a,z)\}/(nK^{n-1})$.
Zero-probability observations contribute zero.

Summing gives the optimal expected error against $\pi$. The finite minimax
theorem allows us to maximize this value over priors to obtain the
smallest worst-case error, proving~\eqref{eq:dual}.
\end{proof}

For a literal linear program, introduce a variable $u_{i,a,z}$ for each
observation. Maximize $(nK^{n-1})^{-1}\sum_{i,a,z}u_{i,a,z}$ over
$\pi\in\Delta(W_S)$ and $u_{i,a,z}\ge0$, subject to
$u_{i,a,z}\le q_0^\pi(i,a,z)$ and
$u_{i,a,z}\le q_1^\pi(i,a,z)$. All these constraints are linear.
The prior is used to describe the value; an optimal finite-game strategy
works for every array and does not need to know a prior.

Equivalently, $\tau(S)$ is the largest possible expected posterior
uncertainty about the held-out label. The characterization requires this
uncertainty to vanish uniformly over finite tuples. The program is finite
once $W_S$ is given, but finding $W_S$ from a representation of $\V$ may
be difficult. The result concerns statistical learnability, not efficient
computation.

One can also define an accuracy-dependent parameter
\[
d_\varepsilon(\V)=
\sup\{n\ge1:\text{there is }S\in\X^n
                       \text{ with }\tau(S)>\varepsilon\},
\]
with the supremum of the empty set taken to be zero. The theorem says that
learning is possible exactly when $d_\varepsilon(\V)<\infty$ for every
$\varepsilon>0$. We do not claim that one fixed value of $\varepsilon$
suffices, or that this parameter has a simpler shattering characterization.

\section{Examples}\label{sec:examples}

\subsection{Direct labels}

Suppose $\V=\{v_h:h\in\Hclass\}$, where $v_h(x,a)=h(x)$ for every
query $a$. The observations at the revealed indices are then ordinary
labels. Our finite game becomes the usual transductive problem of
predicting one held-out label, and learnability is equivalent to finite
VC dimension. This connects the criterion to the classical transductive
approach of \citet{haussler1994predicting}.

\subsection{Independent hidden tables}

We recall a class with two constant hypotheses for which the finite
parameter never decreases. Let $\X=\mathbb N$, $\C=\mathbb Z_4$, and
\[
B=\begin{pmatrix}
0&0&0&0\\
0&0&1&1\\
0&1&1&1\\
0&1&1&1
\end{pmatrix}.
\]
The table has eight zeros and eight ones. Its all-zero first row and first
column ensure that it validates $0$; its complement validates $1$.
For $b\in\{0,1\}$ and maps $\alpha,\beta:\mathbb N\to\mathbb Z_4$,
define
\[
v_{b,\alpha,\beta}(x,p,r)
=b\oplus B_{p+\alpha(x),\,r+\beta(x)},
\qquad
\V_{\mathrm{ind}}=
\{v_{b,\alpha,\beta}:b\in\{0,1\},\
  \alpha,\beta:\mathbb N\to\mathbb Z_4\},
\]
where indices are modulo $4$ and $\oplus$ is addition modulo $2$. Row and column shifts
preserve validity, so the induced hypotheses are exactly $h_b(x)=b$.

Fix a tuple of distinct instances. Choose $b$ fairly, and independently
choose both shifts uniformly at each instance in the tuple. This defines
a prior on $W_S$. Conditional on either value of $b$, every revealed entry
is a fair bit, even when its query is known. The entries observed at
different instances are independent. The full observation therefore gives
no information about $b$, and the held-out target has posterior
probabilities $1/2,1/2$. Proposition~\ref{prop:dual} gives $\tau(S)=1/2$.
There are arbitrarily large distinct tuples, so
$\tau_n(\V_{\mathrm{ind}})=1/2$ for every $n$, although the induced
hypothesis class has VC dimension $1$.

For any $k>4$, the same construction works by fixing a surjection
$\phi:\C\to\mathbb Z_4$ and replacing $p,r$ in the displayed formula by
$\phi(p),\phi(r)$. The extra certificates duplicate rows and columns and
do not reveal additional information.

\subsection{One table shared across instances}

Keep the same base table $B$, but require the shifts to be constant across
instances. Define
\[
\V_{\mathrm{shared}}=
\{v_{b,u,w}(x,p,r)=b\oplus B_{p+u,\,r+w}:
     b\in\{0,1\},\ u,w\in\mathbb Z_4\}.
\]
This class has the same two induced hypotheses as
$\V_{\mathrm{ind}}$, and at each individual instance it permits exactly
the same tables. The difference is that all instances use the same table.

The $n-1$ revealed queries in the finite game now sample entries of that
single table. Once every entry has appeared, its validated label is known.
Predict that label when the table is complete, and otherwise predict a fair
bit. By the union bound, for every tuple $S$,
\[
\tau(S)\le
\frac{K}{2}\left(1-\frac1K\right)^{n-1},
\qquad K=16.
\]
Thus $\tau_n(\V_{\mathrm{shared}})\to0$. The finite pattern families
distinguish these two classes because they retain the relation between
tables at different instances. Neither the induced hypothesis class nor
the set of possible tables at each individual instance makes this
distinction.

\appendix
\section{Passive observations and cooperative learning}\label{app:passive}

We show that uniform random queries preserve learnability. The conversion
from online mistakes to the error of a saved predictor is the standard
online-to-batch argument; see \citet{cesabianchi2004generalization}. Here we
also need to simulate the learner's chosen queries from passive data.

\begin{lemma}\label{lem:passive-equivalence}
The verifier class $\V$ is uniformly iid learnable if and only if
$r_m(\V)\to0$.
\end{lemma}
\begin{proof}
Suppose first that $r_m(\V)\to0$. For each $m\ge1$, choose a passive
learner with worst-case expected error at most $r_m(\V)+1/m$; the
infimum defining $r_m$ need not be attained. At round $t\ge2$, use the
learner for $m=t-1$ on the past observations to predict the current label.
Then query a fresh uniformly random element of $\A$. Since the query
distribution is independent of the prediction and history, the past
observations have exactly the passive iid law. Hence
\[
\E[M_T]\le1+\sum_{m=1}^{T-1}\bigl(r_m(\V)+1/m\bigr)=o(T).
\]

Conversely, suppose an online learner has expected mistake bound $R(T)$
uniformly over $D$ and $v^*$. Fix a virtual horizon $T$. We first simulate
$T$ rounds using an infinite stream of passive observations. At the start
of each virtual round, draw and hold fixed the learner's random seed for
that round. Its proposed prediction and query are then fixed functions of
the candidate instance and the current virtual state. Given a candidate
observation $(X,A,Z)$, expose only
$X$ to the simulated learner and let it propose its prediction and query
$A^*$. Its internal randomness is independent of the recorded $A$.
Accept the round if $A=A^*$, in which case reveal $Z$ and update the
virtual state. Otherwise discard the candidate and all tentative state
changes, retaining the same seed for the next attempt.

Conditional on the virtual state, the candidate instance, and the proposed
prediction and query, acceptance has probability $1/K$. Thus acceptance
does not bias the instance or the learner's proposed choices. The accepted
rounds have exactly the original online learner's law on iid instances.
If $L_T$ is the number of passive observations needed to accept $T$ rounds,
each virtual round has expected waiting time $K$, so $\E[L_T]=KT$.

Let $f_t$ be the saved prediction rule just before the instance of virtual
round $t$ is presented, including that round's fixed seed. Define its population error
by $L(f_t)=\Pr_{X\sim D}(f_t(X)\ne h_{v^*}(X))$. Independence of the
next accepted instance from the virtual history gives
\[
\E[L(f_t)]
=\E\bigl[\ind\{f_t(X_t)\ne h_{v^*}(X_t)\}\bigr].
\]
Consequently, choosing a uniform independent $J\in[T]$ after the ideal
simulation gives $\E[L(f_J)]\le R(T)/T$. Evaluating a saved rule does not
update its state or make a new verifier query.

Now give the simulator only $N$ passive observations. If it finishes $T$
virtual rounds, return the saved rule indexed by $J$; otherwise return an
arbitrary predictor. Couple this learner to the infinite-stream simulation.
On $\{L_T\le N\}$ their outputs agree. Since errors are at most one,
Markov's inequality gives
\[
r_N(\V)
\le\frac{R(T)}{T}+\Pr(L_T>N)
\le\frac{R(T)}{T}+\frac{KT}{N}.
\]
Taking $T=\lfloor\sqrt{N/K}\rfloor$ for sufficiently large $N$ proves
$r_N(\V)\to0$. This coupling avoids conditioning the ideal learner's
error guarantee on the event that the simulation finishes within its
budget.
\end{proof}

\begin{thebibliography}{3}

\bibitem[Asilis et~al.(2024)Asilis, Devic, Dughmi, Sharan, and Teng]{asilis2024transductive}
Julian Asilis, Siddartha Devic, Shaddin Dughmi, Vatsal Sharan, and
Shang-Hua Teng.
\newblock Transductive learning is compact.
\newblock In \emph{Advances in Neural Information Processing Systems},
volume~37, 2024.
\newblock \href{https://proceedings.neurips.cc/paper_files/paper/2024/file/eb7295a8bc613b375726659c2ecd6f14-Paper-Conference.pdf}{Conference paper}.

\bibitem[Cesa-Bianchi et~al.(2004)Cesa-Bianchi, Conconi, and Gentile]{cesabianchi2004generalization}
Nicol\`o Cesa-Bianchi, Alex Conconi, and Claudio Gentile.
\newblock On the generalization ability of on-line learning algorithms.
\newblock \emph{IEEE Transactions on Information Theory}, 50(9):2050--2057,
2004.
\newblock \href{https://cesa-bianchi.di.unimi.it/Pubblicazioni/J20.pdf}{Author's copy}.

\bibitem[Haussler et~al.(1994)Haussler, Littlestone, and Warmuth]{haussler1994predicting}
David Haussler, Nick Littlestone, and Manfred~K. Warmuth.
\newblock Predicting $\{0,1\}$-functions on randomly drawn points.
\newblock \emph{Information and Computation}, 115(2):248--292, 1994.
\newblock \href{https://doi.org/10.1006/inco.1994.1097}{doi:10.1006/inco.1994.1097}.

\end{thebibliography}
\end{document}
